% !PRINTED 69
\weberpage{69}

\begin{center}
\textbf{\S\,20. \gesperrt{The $\vartheta$-Function.}}
\end{center}

The function $\vartheta(v)$ is an \gesperrt{integral} function of $v$ which
satisfies the conditions:

\weq{1}{\begin{aligned}
\vartheta(v + 1) &= e^{-\pi i g_1}\vartheta(v),\\
\vartheta(v + \omega) &= e^{-\pi i(2v+\omega)} e^{-\pi i g_2}\vartheta(v),
\end{aligned}}

\gesperrt{and it is thereby determined up to a factor independent of $v$.}
It is therefore contained among the $t$-functions as a special
case, while on the other hand the general $t$-function can be
expressed by the $\vartheta$-function in the following manner:
\[
t(u) = C\, e^{-\pi i(\lambda u^2 + \mu u)}\, \vartheta\!\left(\frac{u}{\omega_1},\ \frac{\omega_2}{\omega_1}\right),
\]
in which $C, \lambda, \mu$ are arbitrary quantities, either constant or
dependent on $\omega_1, \omega_2$.

The function $\vartheta$ depends further on the numbers $g_1, g_2$
occurring in the characteristic, and if a notation for
this dependence is required, then for $\vartheta(v, \omega)$ we shall set:
\[
\vartheta_{g_1,\, g_2}(v, \omega)
\]
where $g_1, g_2$ may be arbitrary numbers.
According to what has gone before, it suffices to assume them real and between 0
and 2.

Corresponding to the $T$-functions of higher order, we shall
also introduce $\Theta$-functions of the $m$th order, and understand
by these an integral function of $v$ which satisfies the conditions:

\weq{2}{\begin{aligned}
\Theta(v + 1) &= e^{-\pi i g_1}\Theta(v)\\
\Theta(v + \omega) &= e^{-m\pi i(2v+\omega)} e^{-\pi i g_2}\Theta(v).
\end{aligned}}

Here too the characteristic $g_1, g_2$ can be taken up into the
notation:
\[
\Theta_{g_1,\, g_2}(v, \omega).
\]

These functions are contained as special cases among the $T$-functions,
and one can form general $T$-functions in the
following manner:

\weq{3}{C\, e^{-\pi i(\lambda u^2 + \mu u)}\, \Theta\!\left(\frac{u}{\omega_1},\ \frac{\omega_2}{\omega_1}\right).}

There follows therefore from \S\,18, 3. the \gesperrt{fundamental theorem} for
the $\Theta$-functions, which we state as follows:
%CONT>>
